% =========================================================================
% DARCLOS-TEMPO paper draft
%
% Compile: pdflatex paper && bibtex paper && pdflatex paper && pdflatex paper
%          (or simply: make)
%
% This file uses the standard `article` class so it builds with any TeXLive
% installation. For AMT submission, swap the first three lines to
%     \documentclass[journal abbreviation,manuscript]{copernicus}
% (requires the Copernicus package).
% =========================================================================

\documentclass[11pt,a4paper]{article}

\usepackage[utf8]{inputenc}
\usepackage[T1]{fontenc}
\usepackage{lmodern}
\usepackage[a4paper,margin=1in]{geometry}
\usepackage{amsmath,amssymb}
\usepackage{graphicx}
\usepackage{booktabs}
\usepackage{tabularx}
\usepackage{xcolor}
\usepackage{siunitx}
\usepackage{microtype}
\usepackage[round,authoryear]{natbib}
\usepackage{authblk}
\usepackage{caption}
\captionsetup{font=small,labelfont=bf}
\usepackage[hidelinks,colorlinks=false,bookmarks=true]{hyperref}
\usepackage{cleveref}

% --- Custom commands ---------------------------------------------------------

% A grey dashed box used in place of unfinished figures. Replace each
% \placeholderfig call with the matching \includegraphics line when the
% figure is ready.
\newcommand{\placeholderfig}[2][8cm]{%
  \noindent
  \fcolorbox{black!50}{black!4}{%
    \begin{minipage}[c][#1][c]{\dimexpr\linewidth-2\fboxsep-2\fboxrule\relax}
      \centering
      {\sffamily\large\textbf{[ figure placeholder ]}}\\[1em]
      {\itshape\small #2}
    \end{minipage}%
  }%
}

\newcommand{\code}[1]{\texttt{#1}}
\newcommand{\degc}{$^{\circ}$}

% --- Title and authors -------------------------------------------------------

\title{\bfseries DARCLOS\,--\,TEMPO: A cloud shadow detection algorithm adapted
from TROPOMI to the geostationary TEMPO instrument}

\author[1]{M. Pérez-Carrasco}
\author[1]{Q. Zhu}

\affil[1]{[Affiliation TBD]}

\date{Draft, June 2026}

% =========================================================================
\begin{document}
% =========================================================================

\maketitle

\begin{abstract}
\noindent
Cloud shadows are a systematic, currently unmitigated error source in the
trace-gas retrievals of the Tropospheric Emissions: Monitoring of Pollution
(TEMPO) mission. A pixel under the shadow of a nearby cloud is darker than
its clear-sky climatology predicts, biasing the air-mass-factor estimate
used to convert slant to vertical columns of NO$_2$, HCHO, SO$_2$, and O$_3$.
We adapt the DARCLOS cloud-shadow detection algorithm of
\citet{trees2022darclos} to the geostationary UV-VIS spectrometer geometry
of TEMPO. Three structural changes are needed. First, physical cloud heights
are derived from the operational Optical Centroid Pressure (OCP) reported
by the V4 CLDO4 product \citep{wang2025cldo4}, since the corresponding L1B
\code{cloud\_top\_height} field is fill-valued on most pixels. Second, the
directional Lambertian-equivalent reflectance (DLER) climatology used by
DARCLOS for the radiometric shadow confirmation is not available for TEMPO;
in its place, we use the operational effective cloud fraction (ECF) as a
single-wavelength surrogate, because the CLDO4 ECF retrieval already
inverts the measured radiance against the GLER reference and therefore
encodes the same darker-than-climatology signal. Third, viewing and solar
azimuth angles are read from the L1B band group, because the corresponding
L2 geolocation fields are largely fill-valued in V03 and V04. The
implementation is fully vectorized, uses an ECEF-coordinate k-d tree for
batched nearest-pixel lookup, and offers per-granule and per-scan
operational modes; in the latter, L1B and L2 arrays are concatenated along
the mirror-step axis before the algorithm runs, allowing shadows to be
projected across granule boundaries. We present sample results validated
by visual inspection against TEMPO RGB composites on full-day scans of
17--18 September 2025 and discuss the limitations of the single-wavelength
surrogate, the line-based umbra projection, and the bounding-box
georeferencing inherent in the swath-grid output.

\medskip
\noindent\textbf{Keywords:} TEMPO; DARCLOS; cloud shadow detection;
geostationary remote sensing; CLDO4; effective cloud fraction;
air mass factor; trace gas retrievals.
\end{abstract}

% =========================================================================
\section{Introduction}\label{sec:intro}
% =========================================================================

The Tropospheric Emissions: Monitoring of Pollution (TEMPO) mission
\citep{zoogman2017tempo}, launched in April 2023 and operational since
August 2023, is the first geostationary UV-VIS atmospheric composition
spectrometer over North America. From its parking position at 91\degc W,
TEMPO observes NO$_2$, HCHO, SO$_2$, O$_3$, glyoxal, water vapour, and
aerosol indices once per hour at a sub-urban spatial resolution of
$\approx 2.0$\,km north--south by $\approx 4.75$\,km east--west at field
centre. The hourly cadence has no precedent for trace-gas measurements at
this resolution; it is intended to resolve the diurnal cycle of emissions
and atmospheric chemistry over the United States, southern Canada,
Mexico, and the Greater Antilles.

The same hourly cadence over a fixed footprint compounds systematic
errors that correlate with the diurnal evolution of clouds. Cloud shadows
are one such error. A shadow is not a cloud in the line of sight; it is a
region of the surface that receives reduced direct solar illumination
because an adjacent cloud blocks the sun. From the satellite's
perspective the shadowed pixel is anomalously dark, but the darkening
originates at the surface rather than from atmospheric extinction along
the line of sight. The Independent Pixel Approximation (IPA) framework
used in the CLDO4 cloud product
\citep{vasilkov2018cldo, wang2025cldo4} absorbs part of this darkening
into the retrieved Effective Cloud Fraction (ECF), driving it toward
zero, or toward slightly negative values before clipping. A
shadow-contaminated pixel therefore passes the standard NO$_2$ quality
filter of ECF\,$<0.2$ while remaining radiometrically biased. The bias
then propagates into the air-mass-factor (AMF) calculation, which
treats the climatological GLER as the surface reflectance, and from
there into the retrieved trace-gas column.

The problem is not specific to TEMPO. For TROPOMI in low Earth orbit
(7\,$\times$\,3.5\,km$^2$ at nadir), \citet{trees2022darclos} developed
the DARCLOS algorithm, validated against VIIRS true-colour imagery with
$F_1$ scores between 0.84 and 0.95 depending on cloud type and surface.
DARCLOS produces three per-pixel flags: PCSF, a geometric Potential
Cloud Shadow Flag; ACSF, a radiometric Actual Cloud Shadow Flag that
refines PCSF by requiring the scene to be measurably darker than a DLER
climatology at a chosen wavelength; and SCSF, the wavelength-resolved
version of ACSF that produces one flag per DOAS retrieval window. The
DARCLOS paper notes that the algorithm ``can also be applied to other
spectrometers with sufficient spatial resolution''. TEMPO's resolution
satisfies that requirement; the question is what to do about the
ancillary products that DARCLOS uses but TEMPO either does not provide
or provides in a different form.

We identify three differences between TROPOMI and TEMPO that drive the
adaptation:

\begin{enumerate}
  \item TROPOMI uses FRESCO-S physical cloud-top heights directly; the
        TEMPO CLDO4 product reports Optical Centroid Pressure (OCP)
        instead, and the L1B \code{cloud\_top\_height} field is
        fill-valued on most pixels in V03 and V04.
  \item DARCLOS confirms shadows by comparing TOA-derived scene LER
        against a directional LER (DLER) climatology at 13 wavelengths.
        TEMPO operates on the Geometry-dependent LER (GLER) of
        \citet{vasilkov2017gler}, which is provided only at 440 and
        466\,nm in CLDO4 V04.
  \item TEMPO is geostationary. Viewing geometry is fixed per pixel for
        the lifetime of the mission; solar geometry changes continuously
        across the day.
\end{enumerate}

\Cref{sec:background} reviews the relevant elements of DARCLOS and CLDO4.
\Cref{sec:methodology} describes the adaptation: an OCP-to-height
hypsometric conversion (\cref{sec:height}), an ECF-based ACSF surrogate
that relies on the CLDO4 MLER inversion itself (\cref{sec:acsf}), a
vectorized geometric projection with ECEF k-d-tree nearest-pixel lookup
(\cref{sec:projection}), and the two operational modes of the pipeline
(\cref{sec:modes}). \Cref{sec:results} presents sample results on
17--18~September 2025 with visual validation against TEMPO RGB
composites. \Cref{sec:limitations} discusses the theoretical and
implementation limitations of the current approach.
\Cref{sec:conclusions} outlines a prioritised roadmap toward quantitative
validation against GOES-East ABI and the recovery of the polygonal umbra
footprint.

% =========================================================================
\section{Background}\label{sec:background}
% =========================================================================

\subsection{Cloud-shadow geometry}\label{sec:geometry}

A cloud at height $h_c$ above the surface, illuminated by the sun at
solar zenith angle $\theta_0$ and solar azimuth angle $\phi_0$, casts a
shadow whose horizontal displacement from the cloud's nadir point on the
surface is
\begin{equation}
\Delta x_{\rm sh} = -h \tan\theta_0 \sin\phi_0,
\qquad
\Delta y_{\rm sh} = -h \tan\theta_0 \cos\phi_0,
\label{eq:shadow-from-cloud}
\end{equation}
where $h$ is an effective cloud height above the local surface
$h_{\rm sfc}$. A satellite viewing at zenith $\theta$ and azimuth $\phi$
sees the cloud parallax-displaced from its nadir projection by
\begin{equation}
\Delta x_n = h \tan\theta \sin\phi,
\qquad
\Delta y_n = h \tan\theta \cos\phi.
\label{eq:parallax}
\end{equation}
The shadow's ground location relative to the observed pixel centre is
the sum of these two displacements,
\begin{equation}
\Delta x = h \left(\tan\theta\,\sin\phi - \tan\theta_0\,\sin\phi_0\right),
\quad
\Delta y = h \left(\tan\theta\,\cos\phi - \tan\theta_0\,\cos\phi_0\right).
\label{eq:total-displacement}
\end{equation}
\Cref{eq:parallax,eq:total-displacement} reproduce DARCLOS
Eqs.~2--5. The conversion to geodetic offsets uses the local Earth radii
of curvature $M$ and $N$ at the cloud latitude $\phi_c$:
\begin{equation}
\Delta\phi = \frac{\Delta y}{M + h_{\rm sfc}},
\qquad
\Delta\lambda = \frac{\Delta x}{(N + h_{\rm sfc})\cos\phi_c}.
\label{eq:geodetic-offsets}
\end{equation}
At the kilometre scale at which we operate,
$M \approx N \approx 6\,378\,137$\,m and the cosine correction
dominates.

\subsection{The DARCLOS three-flag framework}\label{sec:darclos}

DARCLOS introduces a safety factor $C = 0.5$ that enlarges the
effective height used in the projection,
\begin{equation}
h = (1+C)(h_c - h_{\rm sfc}),
\label{eq:safety-factor}
\end{equation}
which is DARCLOS Eq.~1. The safety factor absorbs the difference between
the cloud centroid (the retrieved quantity) and the geometric cloud top
(the surface that casts the umbra), and accommodates sub-pixel cloud
structure. With the projected shadow on the surface, the PCSF is raised
on every cloud-free pixel inside the triangle defined by the observer
$O$, the cloud-top point $P$, and the shadow tip $Q$.

The ACSF refines PCSF radiometrically. DARCLOS computes the Scene
Normalised Lambertian-Equivalent Reflectance (SCNLER) from the measured
TOA radiance via a Rayleigh-correction lookup table,
\begin{equation}
A_{\rm scene}(\lambda) =
\frac{R^{\rm meas}(\lambda) - R^0(\lambda)}
     {T(\lambda) + s^{\ast}(\lambda)\,T^{\rm meas}(\lambda) - R^0(\lambda)},
\label{eq:scnler}
\end{equation}
and compares it with the directional LER climatology of
\citet{tilstra2017dler}. The contrast
\begin{equation}
\Gamma(\lambda) =
\frac{A_{\rm scene}(\lambda) - A_{\rm DLER}(\lambda)}
     {A_{\rm DLER}(\lambda)} \times 100\,\%
\label{eq:gamma}
\end{equation}
is evaluated at $\lambda_{\max} = \arg\max_\lambda A_{\rm DLER}(\lambda)$,
where shadow detectability is largest. The ACSF is raised when
$\Gamma(\lambda_{\max}) < q = -15\,\%$. The SCSF repeats this check at 13
individual wavelengths in the 328--494\,nm range, producing per-window
flags for the downstream NO$_2$, HCHO, SO$_2$, and O$_3$ retrievals.

\subsection{The TEMPO CLDO4 cloud product}\label{sec:cldo4}

The TEMPO O$_2$--O$_2$ cloud product CLDO4 \citep{wang2025cldo4} inherits
its forward model from the OMI O$_2$--O$_2$ algorithm of
\citet{vasilkov2018cldo}. Under the Independent Pixel Approximation, the
modelled TOA normalised radiance is
\begin{equation}
I_m = I_g (1 - f) + I_c f,
\label{eq:mler}
\end{equation}
which is CLDO4 Eq.~1. Here $I_g$ is the clear-sky radiance computed over
a Lambertian surface whose reflectance equals the GLER, and $I_c$ is the
cloudy radiance for a Lambertian cloud at scene pressure with fixed
albedo $R_c = 0.8$. The effective cloud fraction $f$ (ECF) is retrieved
at 466\,nm by inverting \cref{eq:mler} given the measured $I_m$. The
Optical Centroid Pressure (OCP) is retrieved separately from the
O$_2$--O$_2$ slant column density at 477\,nm.

The ECF inversion is the entry point for our radiometric surrogate. The
inversion uses the GLER as its clear-sky surface reference. A pixel that
is geometrically clear but radiometrically darkened by a shadow has
$I_m < I_g$, and the inversion drives $f$ toward zero, or toward small
negative values before clipping. The CLDO4 ECF therefore encodes
$\Gamma(466\,\text{nm}) < 0$ implicitly, without requiring an explicit
SCNLER--DLER comparison.

CLDO4 also exposes two scene-level diagnostic fields, \code{SceneLER466}
and \code{SceneLER440}, that are central to DARCLOS-style work but are
not used in operational retrievals. \code{SceneLER466} is the
Lambertian-equivalent reflectance that reproduces $I_m$ assuming full
cloud cover. For a clear, unshadowed pixel it equals \code{GLER466}
(the paper's internal consistency check); for a shadowed pixel it falls
below \code{GLER466}. We do not use these fields in the present
implementation but discuss them in \cref{sec:conclusions} as the natural
input for a multi-wavelength refinement of the surrogate.

% =========================================================================
\section{Methodology}\label{sec:methodology}
% =========================================================================

\subsection{Overview}\label{sec:overview}

The adaptation retains the two-stage geometric-then-radiometric structure
of DARCLOS but replaces three components to align with the TEMPO data
products. The three replacements are:

\begin{enumerate}
  \item \textbf{Cloud height} from CLDO4 OCP via the hypsometric equation,
        capped at 16\,km and enlarged by the DARCLOS safety factor
        $C = 0.5$ (\cref{sec:height}).
  \item \textbf{Radiometric refinement} via the CLDO4 ECF itself, used as
        a single-wavelength surrogate for the DARCLOS
        $\Gamma(\lambda_{\max}) < -15\,\%$ criterion. The physical basis
        is that the ECF inversion of \cref{eq:mler} already compares the
        measured radiance with the GLER reference and therefore encodes
        the same darker-than-climatology signal (\cref{sec:acsf}).
  \item \textbf{Vectorised projection} of \cref{eq:total-displacement,eq:geodetic-offsets}
        over all cloud-source pixels in one NumPy call, with a single
        batched \code{cKDTree.query} in ECEF coordinates replacing the
        per-pixel \code{argmin}. The cloud-to-shadow umbra is rasterised
        as a Bresenham line. Angles come from the L1B band group, not
        from L2 geolocation, because the L2 azimuth fields are largely
        fill-valued in V03 and V04 (\cref{sec:angles,sec:projection}).
\end{enumerate}

\Cref{fig:pipeline} shows the pipeline schematically.

\begin{figure}[tb]
\centering
\placeholderfig[7.5cm]{Pipeline schematic. Inputs:
TEMPO L1B (RGB, angles, \code{cloud\_top\_height}) and CLDO4 L2
(\code{cloud\_fraction}, \code{cloud\_pressure}, \code{surface\_pressure},
\code{terrain\_height}, lat/lon). Stages: (a) cloud height; (b) vectorised
DARCLOS projection of \cref{eq:total-displacement,eq:geodetic-offsets};
(c) ECEF k-d-tree nearest-pixel mapping; (d) Bresenham line rasterisation
producing PCSF; (e) ACSF refinement via ECF\,$< \tau_{\rm shadow}$.
Outputs: per-granule or per-scan GeoTIFF plus three shapefiles.}
\caption{DARCLOS-TEMPO pipeline overview.}
\label{fig:pipeline}
\end{figure}

\subsection{Cloud height from optical centroid pressure}\label{sec:height}

The hypsometric equation relates pressure to altitude in an isothermal
atmosphere with constant scale height $H_{\rm scale}$:
\begin{equation}
h_{\rm above\,terrain} =
H_{\rm scale}\,\ln\!\left(\frac{p_{\rm sfc}}{p_{\rm cloud}}\right).
\label{eq:hypsometric}
\end{equation}
We take $H_{\rm scale} = 8500$\,m. The cloud altitude above sea level is
$h_c^{\rm ASL} = h_{\rm above\,terrain} + h_{\rm terrain}$, capped at
$h_{\rm cap} = 16$\,km to bound deep-convection cases in which OCP can
reach very low values. The effective shadow-projection height is then
\begin{equation}
h = (1 + C)\,(h_c^{\rm ASL} - h_{\rm terrain}),
\qquad C = 0.5,
\label{eq:effective-height}
\end{equation}
following the DARCLOS safety-factor convention of \cref{eq:safety-factor}.

The height computation has two branches with explicit priority. Where
the L1B \code{cloud\_top\_height} is finite and above the local terrain,
it is used directly, with the same $h_{\rm cap}$ ceiling and the same
$(1+C)$ multiplier. Where it is fill-valued, which is the case for most
pixels in V03 and V04, the hypsometric expression of
\cref{eq:hypsometric} is used. Where the cloud is retrieved at or below
the surface ($p_{\rm cloud} \geq p_{\rm sfc}$, common for marine stratus
in CLDO4), $h$ is set to zero, which collapses the projection and casts
no shadow.

Source pixels for which neither branch yields a physical height are
counted and reported. If the dropped fraction exceeds a configurable
guard $f_{\rm drop\,max}$ for the granule or the scan, processing fails
loudly rather than silently substituting a default height. In practice
this threshold is hit on scans dominated by snow or by very high SZA,
where the OCP retrieval skips a large fraction of cloud pixels (CLDO4
\code{processing\_quality\_flag} bit 13).

A constant scale height is accurate to about 10\,\% in the troposphere
\citep{wallace2006atmospheric}. The DARCLOS safety factor enlarges the
projection by 50\,\%, which absorbs errors at this level. A
co-located MERRA-2 pressure-altitude profile would refine the height
estimate at high latitudes and in winter; it is not required for
first-order shadow detection and is left for future work
(\cref{sec:conclusions}).

\subsection{ECF as a single-wavelength ACSF surrogate}\label{sec:acsf}

Let $A^{\rm scene}(466\,\text{nm})$ be the CLDO4 SceneLER and
$A^{\rm GLER}(466\,\text{nm})$ be the GLER at the same pixel. The
DARCLOS $\Gamma$ contrast at 466\,nm is
\begin{equation}
\Gamma(466\,\text{nm}) =
\frac{A^{\rm scene}(466\,\text{nm}) - A^{\rm GLER}(466\,\text{nm})}
     {A^{\rm GLER}(466\,\text{nm})}.
\label{eq:gamma-466}
\end{equation}
For a clear, unshadowed pixel, the CLDO4 self-consistency check requires
$A^{\rm scene} \approx A^{\rm GLER}$ \citep{wang2025cldo4}, giving
$\Gamma \approx 0$. For a clear pixel under a shadow,
$A^{\rm scene} < A^{\rm GLER}$ and $\Gamma < 0$. Equivalently, with
$I_m < I_g$ in \cref{eq:mler}, the inverted $f$ goes to zero (or to
slightly negative values pre-clipping in V03).

We use the operational ECF directly as a surrogate for the DARCLOS
$\Gamma(\lambda_{\max})$ criterion:
\begin{equation}
\mathrm{ACSF}(\mathbf{x}) =
\mathrm{PCSF}(\mathbf{x})
\;\wedge\;
\mathrm{ECF}(\mathbf{x}) < \tau_{\rm shadow},
\label{eq:acsf}
\end{equation}
with $\tau_{\rm shadow} = 0.05$ in the default configuration.
Cloud-source pixels are identified by
$\mathrm{ECF} \geq \tau_{\rm cloud} = 0.30$, the threshold recommended
by the TEMPO cloud team for treating a pixel as a coherent shadow
caster.

Two consequences are worth stating. First, the surrogate operates at
one wavelength (466\,nm) only; the wavelength-resolved SCSF is not
produced. Second, in V03, $\tau_{\rm shadow} = 0.05$ overlaps with the
documented $\sim$0.05--0.10 positive ECF bias of clear-sky pixels
\citep{wang2025cldo4}, weakening the surrogate's discrimination. V04
reduces this bias substantially, and $\tau_{\rm shadow} = 0.05$ is a
more meaningful discriminator there. We use V04 as the operational
input in \cref{sec:results}.

\subsection{Source of geometric angles}\label{sec:angles}

The projection of \cref{eq:total-displacement,eq:geodetic-offsets} needs
four angles per pixel: solar zenith, solar azimuth, viewing zenith, and
viewing azimuth. The L2 CLDO4 \code{geolocation} group nominally
contains all four. On the granules examined here, however, the
\code{viewing\_azimuth\_angle} field is entirely fill-valued in V03 and
in V04, and \code{solar\_azimuth\_angle} is $\sim$94\,\% fill-valued. A
literal use of these fields propagates NaN through the projection, lands
the \code{argmin} on grid index $(0,0)$ for every cloud pixel, and
produces a fan of Bresenham lines toward the granule corner. The result
is plausible-looking in RGB overlay because the lines cross enough
genuinely dark pixels for the radiometric AND with the ECF threshold to
latch on. It is not a correct geometric projection.

We instead read all four angles from the L1B \code{band\_290\_490\_nm}
group, where every field is populated. The L1B angles share the
per-pixel identifier with L2 geolocation, so no resampling or coordinate
transformation is needed. If any of the four angles is entirely
fill-valued in the L1B file, the granule fails with an explicit error
rather than producing a wrong projection.

\subsection{Vectorised projection, ECEF k-d tree, and Bresenham
rasterisation}\label{sec:projection}

For each cloud-source pixel $\mathbf{x}_k$, the equations of
\cref{sec:geometry} produce a target shadow location
$\mathbf{y}_k = (\phi_k^{\rm sh}, \lambda_k^{\rm sh})$. The DARCLOS
reference implementation iterates the projection per pixel; we evaluate
\cref{eq:total-displacement,eq:geodetic-offsets} with a single NumPy
call over the entire source-pixel array, since each pixel's projection
depends only on its own position, height, and four angles.

Mapping the projected $(\phi^{\rm sh}, \lambda^{\rm sh})$ back to a
grid index needs a nearest-pixel lookup. A naive \code{argmin} over the
full granule grid is $O(NM)$, where $N$ is the number of source pixels
and $M$ the grid size. At the granule scale
($N \approx 10^4$--$10^5$, $M \approx 2.7\times 10^5$) this runs in
minutes; at the scan scale ($M \approx 3\times 10^6$) it is impractical.
We build a \code{scipy.spatial.cKDTree} over the Earth-Centred
Earth-Fixed (ECEF) coordinates of the grid's valid lat/lon pixels and
query all shadow locations in a single batched call,
$O(N \log M)$. ECEF coordinates avoid the $\cos\phi$ Euclidean-distance
bias of a naive $(\phi, \lambda)$ metric.

The shadow umbra is rasterised as a Bresenham line from the
cloud-source pixel to the projected shadow location. All non-cloud
pixels on the line are added to PCSF. The ACSF refinement
(\cref{eq:acsf}) is a boolean AND with the ECF threshold mask. The
implications of projecting the umbra as a line rather than a polygon
are discussed in \cref{sec:lim-theory}.

\subsection{Operational modes: per-granule and per-scan}\label{sec:modes}

The pipeline ships two CLI drivers with identical configuration
contracts. \code{run\_day.py} processes each TEMPO granule (a
$\sim$6.7-minute slice of one scan) independently and writes one GeoTIFF
and three shapefiles per granule. This is appropriate for granule-level
QA and for spot-checking. \code{run\_scan.py} groups all granules
belonging to one scan (typically 9--11 granules covering a full
east--west pass) and concatenates the L1B and L2 arrays along the
mirror-step axis before the algorithm runs. The merged scan is processed
as a single grid.

Per-scan processing has two advantages over post-hoc mosaicking of
per-granule outputs. First, shadows that fall across granule boundaries
are detected; the per-granule mode cannot project a shadow from a cloud
in granule $n$ onto a candidate pixel in granule $n+1$. Second, the
height-validity guard $f_{\rm drop\,max}$ is averaged over the entire
scan, so a single granule with many invalid OCP retrievals no longer
fails the whole pass. The trade-off is that the bounding-box affine
used to georeference the GeoTIFF spans a larger area, propagating
slightly more swath-grid distortion at the edges (see
\cref{sec:lim-impl}).

\subsection{Implementation and software architecture}\label{sec:software}

The pipeline is written in Python~3 with NumPy, SciPy, rasterio, fiona,
and h5py. We open both L1B and L2 files with \code{h5py} rather than
\code{netCDF4.Dataset}; on some shared-storage configurations the
latter raises \verb|[Errno -101] NetCDF: HDF error| because of HDF5
file-locking, while \code{h5py} reads the same files without issue.
Shapefiles are written directly through \code{fiona} to avoid the
NumPy~2.0 incompatibility in \code{GeoDataFrame.to\_file}, which calls
\code{np.array(geom, copy=False)} on the geometry column.

The code separates the algorithm (\code{shadows.py}, pure NumPy/SciPy,
no I/O) from the I/O helpers (\code{io\_utils.py}) and from the two CLI
drivers (\code{run\_day.py} and \code{run\_scan.py}). Algorithm
parameters --- the safety factor $C$, the cloud and shadow ECF
thresholds, the scale height, the maximum drop fraction --- are exposed
as YAML configuration keys with documented defaults traceable to
\citet{trees2022darclos} or to the CLDO4 ATBD of
\citet{wang2025cldo4}. Source code, configurations, and the methodology
technical report are available on the project repository.

% =========================================================================
\section{Results}\label{sec:results}
% =========================================================================

\subsection{Data}\label{sec:data}

We processed all available TEMPO L1B V04 and CLDO4 L2 V04 granules for
two days of September 2025, the 17th and the 18th. Each day contains
14--15 scans of 9--11 granules each, for a total of approximately 250
granules per day covering CONUS, southern Canada, and northern Mexico.
Granules with fully-fill viewing azimuths in the L1B band group are
excluded; this affected a handful of granules near scan-edge terminator
conditions.

\subsection{Per-granule mode}\label{sec:res-granule}

\Cref{fig:granule-morning} and \cref{fig:granule-afternoon} show two
representative per-granule outputs.

\begin{figure}[tb]
\centering
\placeholderfig[8cm]{Per-granule output for granule S004G05 of
2025-09-18, mid-morning over the central United States. Four panels:
(a)~TEMPO RGB composite, gamma-corrected; (b)~CLDO4 ECF
$\geq 0.30$ mask (cloud-source pixels); (c)~PCSF mask, orange;
(d)~ACSF mask, red. The projected shadow direction is consistent
with a morning sun in the south-east: shadows extend toward the
north-west of the casting clouds.}
\caption{Per-granule output, morning scene with cumulus cover over
the central US.}
\label{fig:granule-morning}
\end{figure}

\begin{figure}[tb]
\centering
\placeholderfig[8cm]{Per-granule output for granule S008G04 of
2025-09-18, mid-afternoon over the Great Plains. A frontal cumulus
band traverses the granule. ACSF polygons follow the band on the
sun-opposite side. Four panels as in \cref{fig:granule-morning}.}
\caption{Per-granule output, afternoon frontal cumulus band over
the Great Plains.}
\label{fig:granule-afternoon}
\end{figure}

Across the September 2025 corpus, ACSF flags account for 4--10\,\% of
cloud-source pixels and 0.5--2\,\% of all granule pixels. The PCSF mask
is 3--5$\times$ larger than ACSF, which reflects the PCSF's design
as a conservative geometric pre-filter: PCSF flags every pixel along the
projection line, and only the truly darkened ones survive the ACSF
threshold.

\subsection{Per-scan mode}\label{sec:res-scan}

\Cref{fig:scan-mosaic} shows a per-scan output covering most of CONUS.

\begin{figure}[tb]
\centering
\placeholderfig[8.5cm]{Per-scan output for scan S004 of 2025-09-18,
covering most of CONUS. Four panels: (a)~merged RGB mosaic from
11~granules; (b)~cloud polygons; (c)~potential shadows (PCSF);
(d)~actual shadows (ACSF). Per-scan polygons that cross former granule
boundaries are highlighted; these would have been split by the
per-granule mode.}
\caption{Per-scan output for a full east--west pass.}
\label{fig:scan-mosaic}
\end{figure}

Compared with a post-hoc mosaic of the same scan's per-granule outputs,
the per-scan output shows three differences. First, the pixel size is
consistent across the entire mosaic; the per-granule mosaic resamples
each granule's affine to a common grid, which leaves visible seams at
granule boundaries. Second, shadow polygons extend across former
granule edges wherever the physical umbra crosses the seam. Third, the
per-scan mode does not fail on scans with high per-granule OCP-skip
rates. On the snow-affected scan S001 of 2025-09-18, which covers
Manitoba and Ontario and has $\sim$14\,\% invalid OCP, the per-scan
pipeline completes successfully; the affected pixels are excluded from
shadow casting only.

\subsection{Visual validation}\label{sec:res-validation}

We inspected ten randomly selected granules from 2025-09-17 and ten
from 2025-09-18, covering scan positions S002--S012. For each granule
we checked three properties against the RGB composite:

\begin{enumerate}
  \item \textbf{Direction.} The ACSF polygons extend away from the
        casting cloud on the side opposite the sun: west of north in
        the morning, east of north in the afternoon.
  \item \textbf{Extent.} Shadow polygons are of plausible length given
        the cloud height implied by OCP and the solar zenith angle.
  \item \textbf{Pairing.} For each cloud band visible in the RGB,
        the algorithm produces an associated ACSF polygon on the
        sun-opposite side.
\end{enumerate}

For 18 of 20 inspected granules all three criteria are met. The two
exceptions are granules over open ocean (Gulf of Mexico and the
Florida Atlantic coast) where ACSF over water does not correspond to
visible shadow darkening in the RGB. This is the open-water
false-positive mode discussed in \cref{sec:lim-theory}.

% =========================================================================
\section{Limitations}\label{sec:limitations}
% =========================================================================

We restrict ourselves to limitations of the algorithm and the
implementation as released. Open scientific questions about the role
of shadow contamination in TEMPO retrievals (which trace gases benefit
most, how the bias propagates through AMF computation, how to weight
the SCSF wavelengths for a given DOAS window) are out of scope.

\subsection{Theoretical and algorithmic limitations}\label{sec:lim-theory}

\paragraph{Open water and intrinsically dark surfaces.}
The GLER over open ocean is 0.03--0.05. The ECF retrieval near coasts
and offshore therefore falls naturally below
$\tau_{\rm shadow} = 0.05$ without any shadow being present. Whenever
the geometric PCSF projection crosses open water, the ACSF surrogate
fires. This is the principal precision risk of the current algorithm
and the cause of both visual-validation failures noted in
\cref{sec:res-validation}.

\paragraph{Bright surfaces: snow, ice, deserts.}
Conversely, over surfaces with GLER\,466~$> 0.5$, even a real shadow
does not depress ECF below 0.05. Shadow detection over bright surfaces
remains an open problem in the operational literature; the surrogate
does not attempt it. DARCLOS proper has the same limitation, but its
SCSF check can recover some detections by inspecting shorter
wavelengths where the relative contrast is larger.

\paragraph{Height bias for deep convection.}
OCP is a radiative centroid, not a geometric top. For optically thick
clouds (deep convection, mature thunderstorms), the geometric top can
be several kilometres above OCP, so the projection underestimates
shadow length for these clouds. The safety factor $C = 0.5$ absorbs
part of the bias but not all of it. A cloud-optical-thickness-dependent
height correction would help but lies outside the scope of the
operational CLDO4 product.

\paragraph{No SZA cutoff.}
The algorithm runs at any SZA. Above $\approx 80\degc$ the shadow
projection becomes geometrically unstable and casts very long shadows
that often fall outside the granule footprint or on the night side of
the terminator. The TEMPO L2 QA convention is SZA\,$<75\degc$; applying
the same cutoff is in the roadmap (\cref{sec:conclusions}).

\paragraph{Single-wavelength radiometric check.}
The DARCLOS SCSF concept, multi-wavelength flagging at the DOAS
retrieval wavelengths, lets downstream NO$_2$, HCHO, SO$_2$, and O$_3$
algorithms filter on a per-window basis. Our surrogate operates at
466\,nm only. Aerosol-vs-shadow discrimination, which DARCLOS achieves
by comparing the spectral slope of $\Gamma$ across wavelengths (aerosol
darkens shorter wavelengths more strongly than longer ones), is not
available.

\paragraph{Line vs polygon umbra projection.}
We project the cloud's shadow as a Bresenham line from the source pixel
centre to the projected shadow tip, not as the full umbra quadrilateral
of the four pixel corners as in DARCLOS. The line captures shadow
length but understates width. Real umbras are 1--2 TEMPO pixels wide
cross-track. An earlier polygon implementation produced sparser, more
fragmented ACSF masks because single-pixel-sized polygons combined
with the strict $\tau_{\rm shadow}$ filter dropped too many candidate
pixels; we kept the line for the current release. A revised polygon
implementation with cluster-level union before rasterisation is the
natural next step.

\paragraph{Constant scale height.}
$H_{\rm scale} = 8500$\,m produces approximately 10\,\% errors at the
tropical versus polar tropopauses. A MERRA-2 pressure-altitude profile
co-located with each granule would close this gap; the relevant
pipeline component is not yet implemented.

\paragraph{Smoke and dust aliased as cloud.}
Wildfire smoke and Saharan dust can be retrieved as ECF\,$\geq 0.5$ by
CLDO4 in extreme events \citep[\S 3.1]{wang2025cldo4}. When the
``cloud'' is in fact an aerosol plume, the algorithm projects a
``shadow'' that is the plume's own contribution to scene darkening
superimposed on its position. The current pipeline cannot tell the two
apart.

\subsection{Implementation limitations}\label{sec:lim-impl}

\paragraph{No quantitative validation.}
Visual inspection against the TEMPO RGB composite is necessary but not
sufficient. The recommended validation reference is GOES-East ABI
true-colour imagery, which is geostationary, time-coincident, and
spatially higher resolution (Band~2 at 0.5\,km). A GOES-East
collocation pipeline is in development.

\paragraph{Azimuth-convention assumption.}
The projection assumes the L1B solar and viewing azimuth fields follow
the geographic convention (from North, clockwise, sun direction). A
sign error in the azimuth would project shadows in the wrong direction,
but the Bresenham line and the radiometric AND would still produce a
visually plausible output. The projected direction is consistent with
the sun in RGB overlays (\cref{sec:res-validation}), but a synthetic
single-cumulus ground-truth test has not been performed.

\paragraph{Bounding-box GeoTIFF georeferencing.}
TEMPO's swath is not rectilinear in lat/lon. The GeoTIFF affine is
derived from \code{rasterio.transform.from\_bounds}, which assumes a
regular lat/lon grid and produces 1--2-pixel-scale geometric distortion
at scan edges. Pixel-accurate georeferencing would require GCP-based
warping using the granule's \code{latitude\_bounds} and
\code{longitude\_bounds} corners. The algorithm's masks are exposed as
shapefiles in part for this reason: vector polygons preserve the
precise pixel-centre coordinates regardless of the raster transform.

\paragraph{Always-overwrite policy.}
Each invocation unconditionally overwrites any existing outputs in the
day or scan output directory. This was deliberate during development,
to make algorithm changes easy to validate. For large backfills,
callers must implement resume logic externally.

\paragraph{Tied source and class thresholds.}
A single threshold $\tau_{\rm cloud} = 0.30$ defines both ``this pixel
casts a shadow'' and ``this pixel is a cloud in the output mask''.
Decoupling them (cast at 0.20, classify at 0.30) would let the algorithm
detect shadows from thin cumulus without inflating the cloud label
exposed to downstream consumers.

\paragraph{Bresenham loop remains scalar.}
Height, projection, and nearest-pixel lookup are all vectorised, but
the Bresenham line drawing is a per-source-pixel Python loop. For
$\sim 10^5$ source pixels per granule this is acceptable (a few
seconds); it would be the next optimisation target if the pipeline
were applied to a long-term TEMPO archive.

% =========================================================================
\section{Conclusions and future work}\label{sec:conclusions}
% =========================================================================

We adapted the DARCLOS cloud-shadow detection algorithm
\citep{trees2022darclos} to the geostationary UV-VIS TEMPO instrument.
The adaptation retains DARCLOS's two-stage geometric-then-radiometric
structure but replaces three components to align with the TEMPO data
products. Cloud height is derived from CLDO4 Optical Centroid Pressure
through a hypsometric conversion with the DARCLOS safety factor
preserved. The multi-wavelength SCNLER-vs-DLER contrast is replaced by
a single-wavelength surrogate that uses the operational CLDO4 ECF
retrieval, which already encodes the DARCLOS ``darker than the GLER
climatology'' signal through its MLER inversion. The projection geometry
is vectorised over all cloud-source pixels and accelerated with a
batched ECEF k-d-tree nearest-pixel lookup.

The pipeline runs in two operational modes, per-granule and per-scan
(with L1B/L2 array concatenation before the algorithm runs), and
produces georeferenced GeoTIFF and shapefile outputs that can be
consumed directly by trace-gas retrieval filters or by
machine-learning-based shadow segmentation pipelines. Visual inspection
of CONUS coverage on 17--18~September 2025 shows that the algorithm
correctly identifies the direction, position, and extent of cloud
shadows for resolved cumulus and frontal stratus, with the documented
failure modes over open water, snow, and aerosol plumes.

We do not claim quantitative performance figures because no independent
ground truth has been used. The first next step is GOES-East ABI
validation: the same approach used by \citet{trees2022darclos} against
VIIRS, applied to a representative sample of TEMPO scans to yield
precision, recall, and $F_1$ stratified by cloud type and surface.
Beyond validation, ranked by expected impact:

\begin{enumerate}
  \item \textbf{Explicit quality masking.} SZA\,$<75\degc$,
        \code{snow\_ice\_fraction}\,$\leq 0.5$, and CLDO4
        \code{processing\_quality\_flag} error bits, applied before the
        algorithm runs.
  \item \textbf{Brightness gate for ACSF.} Require GLER\,466~$> 0.05$
        before firing ACSF. The surface should be bright enough that
        a shadow is meaningful. This is the cheapest defence against
        open-water false positives.
  \item \textbf{Multi-wavelength surrogate.} Add an analogous SceneLER\,440
        vs GLER\,440 check, requiring the contrast at both 440 and
        466\,nm. This would discriminate aerosol darkening (steeper
        spectral slope) from shadow darkening (flat slope).
  \item \textbf{Polygonal umbra projection.} Replace the Bresenham line
        with the four-corner umbra polygon, with cluster-level union
        before rasterisation so that adjacent cloud-source pixels
        combine into one continuous polygon.
  \item \textbf{MERRA-2-based height.} Replace the constant scale height
        with a co-located MERRA-2 pressure-altitude profile.
  \item \textbf{GCP-based GeoTIFF georeferencing.} Use the
        \code{latitude\_bounds} and \code{longitude\_bounds} corners
        directly through GDAL warp.
\end{enumerate}

In the longer term, the rigorous path is to replace the
single-wavelength ECF surrogate with a full SCNLER computation from
L1B radiances and a Rayleigh-correction LUT, evaluated at all 13
DARCLOS wavelengths. That depends on operational LUT infrastructure
that is impractical for individual research groups but feasible for
the TEMPO SDPC.

% =========================================================================
\section*{Code and data availability}\addcontentsline{toc}{section}{Code and data availability}
% =========================================================================

The DARCLOS-TEMPO source code, configuration files, the methodology
technical report, and example outputs are available at the project
repository [URL to be inserted]. TEMPO V04 L1B and CLDO4 V04 L2 data
are available through the NASA Atmospheric Science Data Center (ASDC,
\url{https://asdc.larc.nasa.gov}) and the TEMPO Science Data Processing
Center.

% =========================================================================
\section*{Author contributions}\addcontentsline{toc}{section}{Author contributions}
% =========================================================================

[Author contributions to be inserted before submission, following the
CRediT taxonomy.]

% =========================================================================
\section*{Competing interests}\addcontentsline{toc}{section}{Competing interests}
% =========================================================================

The authors declare no competing interests.

% =========================================================================
\section*{Acknowledgements}\addcontentsline{toc}{section}{Acknowledgements}
% =========================================================================

We thank the TEMPO Science Team and the CLDO4 algorithm developers at
the Smithsonian Astrophysical Observatory and the Science Data
Processing Center for making the operational L1B and L2 products
publicly available, and for discussions on the OCP retrieval and on
the SceneLER--GLER consistency check that motivates our ACSF
surrogate.

% =========================================================================
% References
% =========================================================================

\bibliographystyle{abbrvnat}
\bibliography{references}

% =========================================================================
\appendix
% =========================================================================

\section{Algorithm parameters}\label{app:params}

\begin{table}[hb]
\centering
\caption{DARCLOS-TEMPO algorithm parameters as released.}
\label{tab:params}
\begin{tabular}{@{}llll@{}}
\toprule
Parameter & Symbol & Value & Source / rationale \\
\midrule
Cloud-source ECF threshold        & $\tau_{\rm cloud}$    & 0.30        & TEMPO cloud team \\
Shadow ECF threshold (ACSF)       & $\tau_{\rm shadow}$   & 0.05        & DARCLOS-TEMPO; V04-tuned \\
Safety factor                     & $C$                   & 0.5         & \citet{trees2022darclos}, Eq.~1 \\
Hypsometric scale height          & $H_{\rm scale}$       & 8500\,m     & Standard troposphere \\
Cloud altitude cap                & $h_{\rm cap}$         & 16\,000\,m  & Tropopause guard \\
Max drop fraction (per granule/scan) & $f_{\rm drop\,max}$ & 0.10--0.30 & Empirically tuned per V0x \\
Cloud Lambertian albedo (in CLDO4)& $R_c$                 & 0.8         & CLDO4 ATBD \citep{wang2025cldo4} \\
Earth radius ($M$, $N$)           & ---                   & 6\,378\,137\,m & WGS-84 equatorial \\
\bottomrule
\end{tabular}
\end{table}

\section{Output file conventions}\label{app:outputs}

Per-granule outputs (\code{run\_day.py}) land in
\code{\{output\_root\}/\{YYYY\}/\{MM\}/\{DD\}/}:
\begin{verbatim}
rgb_<L1_basename>.tif
<L1_basename>_clouds.shp                (+ .dbf, .shx, .prj, .cpg)
<L1_basename>_potential_shadows.shp
<L1_basename>_actual_shadows.shp
\end{verbatim}

Per-scan outputs (\code{run\_scan.py}) land in
\code{\{output\_root\}\_scan/\{YYYY\}/\{MM\}/\{DD\}/}:
\begin{verbatim}
rgb_S<scan>.tif
S<scan>_clouds.shp
S<scan>_potential_shadows.shp
S<scan>_actual_shadows.shp
\end{verbatim}

Shapefiles use EPSG:4326. The GeoTIFF affine transform is derived from
the bounding box of the granule (or scan) latitude/longitude corners via
\code{rasterio.transform.from\_bounds}. Per-feature attributes include
\code{type}, \code{name}, \code{value}, \code{source} (granule or scan
identifier), and \code{area\_deg2}.

\end{document}
